\subsection{Tight convergence of measures and compact convergence of functions}

PROPOSITION 12. --- Let T be a completely regular space, and let B be the unit ball of the normed space \(\mathcal{C}^{b}(\mathrm{T};\mathbf{C})\) . Let I be a linear form on \(\mathcal{C}^{b}(\mathrm{T};\mathbf{C})\) . In order that there exist a bounded complex measure \(\theta\) on T such that \(\mathrm{I}(f)=\theta(f)\) for all \(f\in\mathcal{C}^{b}(\mathrm{T};\mathbf{C})\) , it is necessary and sufficient that the restriction of I to B be continuous for the topology of compact convergence. The measure \(\theta\) is then unique.

Let us show that the condition in the statement is necessary. Let \(\theta\) be a bounded complex measure on \(\mathbf{T}\) , \(\varepsilon\) a number \(>0\) , and \(\mathbf{K}\) a compact subset of \(\mathbf{T}\) such that \(|\theta|^{\bullet}(\mathbf{T} - \mathbf{K}) < \varepsilon\) . Let \(f \in \mathbf{B}\) ; we denote by \(\mathbf{U}\) the neighborhood of \(f\) in \(\mathbf{B}\) for the topology of compact convergence, formed by the functions \(g \in \mathbf{B}\) such that \(\sup_{x \in \mathbf{K}} |g(x) - f(x)| \leqslant \varepsilon\) . Then, for every \(g \in \mathbf{U}\) ,

\[
| \theta (g) - \theta (f) | \leqslant \int_ {\mathrm{T}} | g - f | d | \theta | \leqslant \varepsilon | \theta | ^ {\bullet} (\mathrm{K}) + 2 | \theta | ^ {\bullet} (\mathrm{T} - \mathrm{K}) \leqslant (\| \theta \| + 2) \varepsilon ,
\]

because \(|g - f|\) is bounded above by \(\varepsilon\) on \(K\) and by 2 on \(T - K\) .

Conversely, consider a linear form I on \(\mathcal{C}^{b}(\mathrm{T};\mathbf{C})\) whose restriction to B is continuous for the topology of compact convergence. Then, for every number \(\varepsilon > 0\) , there exist a number a \textgreater{} 0 and a compact subset K of T such that the relations \(f \in B\) , \(\sup_{x \in K} |f(x)| \leqslant a\) imply \(|\mathrm{I}(f)| \leqslant \varepsilon\) . Prop. 5 of No.~2 then implies the existence of a unique bounded complex measure \(\theta\) such that \(\mathrm{I}(f) = \theta(f)\) for all \(f \in \mathcal{C}^{b}(\mathrm{T};\mathbf{C})\) .

PROPOSITION 13. --- Let T be a locally compact space, and H a bounded subset of the normed space \(\mathcal{C}^{b}(\mathrm{T};\mathbf{C})\) . The mapping \((\mu,f)\mapsto\mu(f)\) of \(\mathcal{M}_{+}^{b}(\mathrm{T})\times\mathrm{H}\) into C is then continuous, when \(\mathcal{M}_{+}^{b}(\mathrm{T})\) is equipped with the tight topology, and H with the topology of compact convergence.

Let \(\mu \in \mathcal{M}_{+}^{b}(\mathrm{T})\) , \(f \in \mathrm{H}\) , and let \(\mathbf{M}\) be a real number such that \(\| \mu \| < \mathbf{M}\) , and \(|g| \leqslant \mathbf{M}\) for all \(g \in \mathbf{H}\) . Let \(\varepsilon\) be a number \(>0\) and choose a compact subset \(\mathbf{K}\) of \(\mathbf{T}\) such that \(\mu^{\bullet}(\mathbf{T} - \mathbf{K}) < \varepsilon\) , then a relatively compact open neighborhood \(\mathbf{S}\) of \(\mathbf{K}\) . The set \(\mathbf{U}\) of measures \(\lambda \in \mathcal{M}_{+}^{b}(\mathbf{T})\) satisfying the inequalities

\[
\lambda^ {\bullet} (\mathrm{T}) <   \mathrm{M}, \quad \lambda^ {\bullet} (\mathrm{T} - \mathrm{S}) <   \varepsilon , \quad | \lambda (f) - \mu (f) | <   \varepsilon
\]

is then a neighborhood of \(\mu\) in \(\mathcal{M}_+^b (\mathrm{T})\) (No.~3, Prop. 6). In addition, let V be the neighborhood of \(f\) in H consisting of the functions \(g\in \mathbf{H}\) such that

\[
\sup _ {x \in S} | g (x) - f (x) | <   \varepsilon .
\]

Let \(\lambda \in \mathbf{U}\) and \(g\in \mathbf{V}\) ; since the function \(|g - f|\) is bounded above by \(\varepsilon\) in \(\mathbf{S}\) , and by 2M in \(\mathbf{T} - \mathbf{S}\) , we have

\[
\left| \lambda (g) - \lambda (f) \right| \leqslant \int_ {\mathrm{T}} | g - f | d \lambda \leqslant \varepsilon \lambda^ {\bullet} (\mathrm{S}) + 2 \mathrm{M} \lambda^ {\bullet} (\mathrm{T} - \mathrm{S}) \leqslant 3 \mathrm{M} \varepsilon ,
\]

from which one deduces

\[
\left| \lambda (g) - \mu (f) \right| \leqslant \left| \lambda (g) - \lambda (f) \right| + \left| \lambda (f) - \mu (f) \right| \leqslant (3 M + 1) \varepsilon .
\]

This proves the continuity of the mapping \((\lambda, g) \mapsto \lambda(g)\) at the point \((\mu, f)\) of \(\mathcal{M}_+^b(\mathrm{T}) \times \mathrm{H}\) .

Remark. --- Let T be a completely regular space, M a subset of \(\mathcal{M}^{b}(\mathrm{T})\) that satisfies Prokhorov's condition, H a bounded subset of \(\mathcal{C}^{b}(\mathrm{T})\) . An argument very close to the one just made may be used to prove that the mapping \((\lambda,g)\mapsto\lambda(g)\) of \(M\times H\) into C is continuous when M is equipped with the tight topology and H with the topology of compact convergence.

COROLLARY. --- Let T be a completely regular space, X a topological space, and f a complex-valued function defined on \(T \times X\) , continuous and bounded. For every bounded measure \(\mu\) on T, let \(F_{\mu}\) be the function on X defined by \(\mathrm{F}_{\mu}(x) = \int_{\mathrm{T}} f(t, x) d\mu(t)\) for all \(x \in X\) .

\begin{enumerate}
\def\labelenumi{\alph{enumi})}
\item
  The function \(\mathrm{F}_{\mu}\) is continuous and bounded for every bounded measure \(\mu\) .
\item
  Suppose that T is locally compact. The mapping \(\mu\mapsto F_{\mu}\) of \(\mathcal{M}_{+}^{b}(T)\) into \(\mathcal{C}^{b}(X;C)\) is then continuous, if \(\mathcal{M}_{+}^{b}(T)\) is equipped with the tight topology, and \(\mathcal{C}^{b}(X;C)\) with the topology of compact convergence.
\end{enumerate}

For every \(x \in X\) , denote by \(f_{x}\) the continuous and bounded function \(t \mapsto f(t, x)\) on T; the mapping \(x \mapsto f_{x}\) of X into \(\mathcal{C}^{b}(\mathrm{T}; \mathbf{C})\) has bounded image, and it is continuous if \(\mathcal{C}^{b}(\mathrm{T}; \mathbf{C})\) is equipped with the topology of compact convergence (GT, X, §3, No.~4, Th. 3). Since \(\mathrm{F}_{\mu}(x) = \mu(f_{x})\) , the function \(F_{\mu}\) is continuous by Prop. 12. Suppose T is locally compact; Prop. 13 shows that the mapping \((\mu, x) \mapsto \mathrm{F}_{\mu}(x)\) of \(\mathcal{M}_{+}^{b}(\mathrm{T}) \times \mathrm{X}\) into C is continuous; the assertion b) follows from this (loc. cit.).

